% !PRINTED 729
\weberpage{729}
%<<CONT
imprimitive Zahlen) enthalten und durch die Gleichung $\mathfrak{E}(r)
= \mathfrak{E}(-r)$ charakterisirt.

Wir wollen nun unter einer \gesperrt{Normaleinheit} des Körpers $H_m$
eine Einheit verstehen, die nicht $= \pm 1$ ist, aber der Gleichung
\weq{3}{\mathfrak{E}(r)\mathfrak{E}(-r) = \pm 1}
genügt.

Es ist klar, dass eine Einheit des Körpers $H_\mu$ dieser Bedingung
 jedenfalls nicht genügen kann. Aber nicht jede Einheit
in $H_m$, die dem Körper $H_\mu$ nicht angehört, wird Normaleinheit
sein $^{1)}$.

Die Einheiten $\tau_\beta$ [\S.\,199, (5)] gehören aber zu den Normaleinheiten
. Es ist nämlich
\weq{4}{\tau_\beta = \mathrm{tg}\,\frac{5^\beta\pi}{m} = \frac{1-r_\beta}{1+r_\beta}r_\beta^{\nu},}
und wenn darin die Substitution $(r,-r) = (r_\beta, r_{\beta+\nu'})$ gemacht
wird, so geht $\tau_\beta$ in
\weq{5}{\tau_{\beta+\nu'} = -\frac{1}{\tau_\beta}}
über, was der Relation (3) entspricht.

Ein System von $\nu'$ Normaleinheiten
\[
\mathfrak{E}_0(r), \mathfrak{E}_1(r), \ldots, \mathfrak{E}_{\nu'-1}(r)
\]
soll ein \gesperrt{unabhängiges System} heissen, wenn die Determinante
\weq{6}{\pm\Sigma\pm\log|\mathfrak{E}_0(r_0)|\log|\mathfrak{E}_1(r_1)|\ldots\log|\mathfrak{E}_{\nu'-1}(r_{\nu'-1})|
= L(\mathfrak{E}_0, \mathfrak{E}_1, \ldots, \mathfrak{E}_{\nu'-1})}
einen von Null verschiedenen Werth hat.

Den absoluten Werth dieser Determinante, $L$, nennen wir
auch hier den \gesperrt{Regulator} des Systems $\mathfrak{E}_0, \mathfrak{E}_1, \ldots, \mathfrak{E}_{\nu'-1}$.

Die Einheiten $\tau_0, \tau_1, \ldots, \tau_{\nu'-1}$ bilden ein unabhängiges System
normaler Einheiten, denn ihr Regulator
\[
L(\tau_0, \tau_1, \ldots, \tau_{\nu'-1})
\]
ist eben die im vorigen Paragraphen definirte Determinante $T$
[\S.\,199, (14), (15)]. Und dass diese nicht Null sein kann, ergiebt
sich unmittelbar daraus, dass sie in dem Ausdrucke für die
%CONT>>

\par\noindent{\footnotesize $^{1)}$ Der Verfasser hat in der oben erwähnten Arbeit (Acta
mathematica, Bd.\,8) die Normaleinheiten als \glqq primitive
Einheiten\grqq{} bezeichnet. Dieser Name ist hier verworfen, weil er
zu dem Irrthum Veranlassung geben könnte, als ob jede Einheit
des Körpers $H_m$, die nicht zugleich dem Körper $H_\mu$ angehört,
dazu zu rechnen sei.}
